\section{Магнитное кольцо: исходная система}
\label{sec:mag-base}

\subsection{Обозначения}

Индекс $s$ отмечает синхронную частицу: импульс $p_s$, скорость $v_s=\beta_s c$, лоренц-фактор $\gamma_s$, энергия $E_s=\gamma_s mc^2$, длина орбиты $C_s$, период обращения $T_s=C_s/v_s$.

Отклонения импульса и энергии частицы:
\begin{equation}
\label{eq:delta-def}
\delta=\frac{p-p_s}{p_s},
\qquad
w=\frac{E-E_s}{\beta_s^2E_s}.
\end{equation}
Из $\mathrm{d}E/\mathrm{d}p=v$ и $\mathrm{d}^2E/\mathrm{d}p^2=c^2/(\gamma^2E)$
\begin{equation}
\label{eq:w-series}
w=\delta+\frac{\delta^2}{2\gamma_s^2}+O(\delta^3),
\qquad
\delta=w-\frac{w^2}{2\gamma_s^2}+O(w^3).
\end{equation}

Фаза $\varphi$ --- фаза ВЧ-поля в момент прохождения частицей зазора, отсчитанная от синхронной частицы; она растёт, когда частица опаздывает. Кратность $h$, частота $\omega_\rf=2\pi h/T_s$, напряжение зазора $V\sin(\varphi_s+\varphi)$, заряд $q$. Номер оборота $k$, время $t=kT_s$. Относительное запаздывание за оборот
\begin{equation}
\label{eq:tau-def}
\tau=\frac{T-T_s}{T_s},
\end{equation}
где $T$ --- период обращения частицы. Рассматривается одна частица, её траектория лежит внутри сепаратрисы.

\subsection{Запаздывание за оборот}

В магнитном поле скорость вдоль траектории постоянна, поэтому $T=C/v$ и точно
\begin{equation}
\label{eq:T-ratio}
\frac{T}{T_s}=\frac{C}{C_s}\,\frac{v_s}{v}.
\end{equation}
Длина замкнутой орбиты записывается рядом
\begin{equation}
\label{eq:C}
\frac{C}{C_s}=1+\alpha_0\delta+\alpha_1\delta^2+O(\delta^3).
\end{equation}
Из $\beta\gamma=\beta_s\gamma_s(1+\delta)$
\begin{equation}
\label{eq:vsv}
\frac{v_s}{v}=\frac{\gamma}{\gamma_s(1+\delta)}
=1-\frac{\delta}{\gamma_s^2}+\Bigl(1+\frac{\beta_s^2}{2}\Bigr)\frac{\delta^2}{\gamma_s^2}+O(\delta^3).
\end{equation}
Произведение \eqref{eq:C} и \eqref{eq:vsv} даёт
\begin{equation}
\label{eq:tau}
\tau=\eta_0\delta+\eta_1\delta^2,
\qquad
\eta_0=\alpha_0-\frac{1}{\gamma_s^2},
\qquad
\eta_1=\alpha_1-\frac{\alpha_0}{\gamma_s^2}+\Bigl(1+\frac{\beta_s^2}{2}\Bigr)\frac{1}{\gamma_s^2}.
\end{equation}
Через $w$ с помощью \eqref{eq:w-series}
\begin{equation}
\label{eq:tau-w}
\tau=\eta_0w+\tilde\eta_1w^2,
\qquad
\tilde\eta_1=\eta_1-\frac{\eta_0}{2\gamma_s^2}.
\end{equation}

\subsection{Система продольных уравнений}

Ускорения нет: $\varphi_s=0$ или $\pi$, и прирост энергии в зазоре $qV\sin(\varphi_s+\varphi)=qV\cos\varphi_s\sin\varphi$. За оборот
\begin{equation}
\label{eq:map}
\varphi_{k+1}-\varphi_k=2\pi h\,\tau_k,
\qquad
w_{k+1}-w_k=\frac{qV}{\beta_s^2E_s}\cos\varphi_s\sin\varphi_k .
\end{equation}
Оба равенства точные: опоздание на $T-T_s$ сдвигает фазу на $\omega_\rf(T-T_s)=2\pi h\tau$, а зазор меняет энергию, то есть $w$, на заданную величину. Деление на $T_s$ даёт систему
\begin{equation}
\label{eq:smooth}
\frac{\mathrm{d}\varphi}{\mathrm{d}t}=\omega_\rf\,\tau,
\qquad
\frac{\mathrm{d}w}{\mathrm{d}t}=K\cos\varphi_s\sin\varphi,
\qquad
K=\frac{qV}{\beta_s^2E_sT_s}.
\end{equation}
Первое уравнение описывает изменение фазы из-за запаздывания, второе --- изменение энергии в зазоре.

Устойчивой точкой $\varphi=w=0$ является
\begin{equation}
\label{eq:phis}
\cos\varphi_s=-\operatorname{sgn}\eta_0 ,
\end{equation}
малые колебания около неё имеют частоту
\begin{equation}
\label{eq:ws}
\omega_s^2=K\omega_\rf|\eta_0|,
\qquad
K\omega_\rf\cos\varphi_s=-\frac{\omega_s^2}{\eta_0}.
\end{equation}

В этой главе $\sin\varphi\approx\varphi$, и с \eqref{eq:tau-w} система принимает вид
\begin{equation}
\label{eq:system}
\frac{\mathrm{d}\varphi}{\mathrm{d}t}=\omega_\rf\bigl(\eta_0w+\tilde\eta_1w^2\bigr),
\qquad
\frac{\mathrm{d}w}{\mathrm{d}t}=K\cos\varphi_s\,\varphi.
\end{equation}

\subsection{Решение}
\label{sec:solution}

Дифференцирование второго уравнения \eqref{eq:system}, подстановка первого и \eqref{eq:ws} дают
\begin{equation}
\label{eq:ode}
\ddot w+\omega_s^2w=-\frac{\omega_s^2\tilde\eta_1}{\eta_0}\,w^2.
\end{equation}
Его решение (приложение~\ref{app:bm}):
\begin{align}
w(t)&=w_\eq+A\cos\psi+\frac{\tilde\eta_1A^2}{6\eta_0}\cos2\psi,
\qquad
w_\eq=-\frac{\tilde\eta_1A^2}{2\eta_0},
\qquad
\psi=\Omega t,
\label{eq:w-sol}\\
\varphi(t)&=\frac{\omega_\rf}{\Omega}\Bigl(\eta_0A\sin\psi+\frac{\tilde\eta_1A^2}{3}\sin2\psi\Bigr),
\label{eq:phi}\\
\frac{\Omega}{\omega_s}&=1-\frac{5}{12}\Bigl(\frac{\tilde\eta_1A}{\eta_0}\Bigr)^2.
\label{eq:Omega}
\end{align}
Амплитуда $A=(w_{\max}-w_{\min})/2$: при $\psi=0$ и $\psi=\pi$ вторая гармоника равна одной и той же величине и полуразмах не меняет. $\Omega$ --- частота синхротронных колебаний, то есть частота обхода траектории на фазовой плоскости $(\varphi,w)$.

\subsection{Усреднение и равновесный сдвиг}

$\langle f\rangle$ --- среднее по оборотам за один синхротронный период. Равновесные сдвиги энергии и импульса:
\begin{equation}
\label{eq:deq-def}
w_\eq=\langle w\rangle,
\qquad
\delta_\eq=\langle\delta\rangle .
\end{equation}
На замкнутой траектории фаза возвращается к прежнему значению, поэтому из первого уравнения \eqref{eq:smooth} точно
\begin{equation}
\label{eq:tau-avg}
\langle\tau\rangle=0 :
\end{equation}
средний период обращения частицы равен $T_s$. С \eqref{eq:tau} и \eqref{eq:tau-w}
\begin{equation}
\label{eq:deq-avg}
\delta_\eq=-\frac{\eta_1}{\eta_0}\langle\delta^2\rangle,
\qquad
w_\eq=-\frac{\tilde\eta_1}{\eta_0}\langle w^2\rangle .
\end{equation}
На решении \eqref{eq:w-sol} $\langle w^2\rangle=\langle\delta^2\rangle=A^2/2+O(A^3)$, откуда
\begin{equation}
\label{eq:deq}
\delta_\eq=-\frac{\eta_1A^2}{2\eta_0},
\qquad
w_\eq=\delta_\eq+\frac{A^2}{4\gamma_s^2}=-\frac{\tilde\eta_1A^2}{2\eta_0},
\qquad
\Delta E_\eq=\beta_s^2E_s\,w_\eq .
\end{equation}
Равенство $w_\eq=\delta_\eq+\langle\delta^2\rangle/(2\gamma_s^2)$ следует из \eqref{eq:w-series}.

Условия применимости:
\begin{equation}
\label{eq:valid}
\Bigl|\frac{\eta_1A}{\eta_0}\Bigr|\ll1,
\qquad
\hat\varphi=\frac{\omega_\rf|\eta_0|A}{\omega_s}\ll1,
\end{equation}
где $\hat\varphi$ --- полуразмах фазы.
